\documentclass[11pt,a4paper]{article}
\usepackage[margin=2.2cm]{geometry}
\usepackage{booktabs}
\usepackage{tabularx}
\usepackage{longtable}
\usepackage{xcolor}
\usepackage{graphicx}
\usepackage{array}
\usepackage{float}
\usepackage[colorlinks=true,linkcolor=blue!50!black,urlcolor=blue!50!black]{hyperref}
\usepackage{caption}

\newcommand{\PASS}{\textcolor{green!50!black}{\textbf{PASS}}}
\newcommand{\FAIL}{\textcolor{red!70!black}{\textbf{FAIL}}}
\newcommand{\OPEN}{\textcolor{orange!80!black}{\textbf{OPEN}}}
\newcommand{\um}{\,$\mu$m}
\newcommand{\umsq}{\,$\mu$m$^2$}
\setlength{\parskip}{4pt}
\setlength{\parindent}{0pt}
\renewcommand{\arraystretch}{1.15}

\title{\textbf{D05 RF Energy Harvester --- Module and Design-Rule Report}\\[4pt]
\large GF180MCU, Chipathon 2026 project slot D05 (variant EH)}
\author{Gelochip / IRIS Harvester}
\date{\today}

\begin{document}
\maketitle

\begin{abstract}
\noindent
This report documents every module built for the D05 tapeout, the design rules each module
must satisfy, whether it satisfies them, and the measured outputs. All results were produced
by the foundry's own GF180MCU KLayout sign-off deck (\texttt{run\_drc.py}), \texttt{netgen}
LVS, \texttt{magic} parasitic extraction and \texttt{ngspice-46}, run on the exact submitted
file \texttt{D05\_Gelochip.gds} and the exact top cell \texttt{D05\_Gelochip}.
\end{abstract}

\section{Scope and submitted files}

\begin{tabularx}{\textwidth}{@{}lX@{}}
\toprule
\textbf{File} & \textbf{Role} \\
\midrule
\texttt{gds/D05\_Gelochip.gds} & \textbf{The tapeout file.} Hierarchical, 163 cells, dummy fill applied, bounding box exactly $(0,0)$--$(550,550)$\um. \\
\texttt{D05\_Gelochip\_nofill.gds} & Same topology without fill. Review aid only --- it fails density. \\
\texttt{D05\_Gelochip\_FET.gds} & Capacitor-free view used for LVS and PEX (see Sec.~\ref{sec:method}). \\
\texttt{D05\_Gelochip\_source.spice} & LVS source netlist. \\
\texttt{RFEH\_PEX\_pex.spice} & Extracted netlist: 2293 resistors, 439 capacitors. \\
\bottomrule
\end{tabularx}

\section{Module catalogue}

Nine modules were built and validated individually. The five transistor blocks are taken
verbatim from the signed-off design; the capacitor modules are new (Sec.~\ref{sec:mim}).

\begin{longtable}{@{}llrrl@{}}
\toprule
\textbf{Module} & \textbf{Function} & \textbf{Size (\um)} & \textbf{Value} & \textbf{Status} \\
\midrule
\endhead
\texttt{unit\_diode}   & guarded diode-connected NMOS      & $10.8\times9.1$    & ---        & \PASS \\
\texttt{rectifier}     & cross-coupled full-wave rectifier & $17.2\times32.9$   & ---        & \PASS \\
\texttt{startup\_bias} & injection-assisted startup bias   & $19.8\times30.9$   & ---        & \PASS \\
\texttt{pump\_ladder}  & 3-stage Dickson charge pump       & $48.3\times14.1$   & ---        & \PASS \\
\texttt{load\_diode}   & MOS-diode output load             & $37.9\times6.1$    & ---        & \PASS \\
\texttt{mimb\_unit}    & \textbf{MIM-B capacitor unit}     & $29.6\times30.1$   & 0.784\,pF  & \PASS \\
\texttt{cap\_inject}   & injection capacitor bank          & $25.7\times19.4$   & 0.300\,pF  & \PASS \\
\texttt{cap\_stage}    & pump stage capacitor bank         & $32.9\times54.5$   & 1.200\,pF  & \PASS \\
\texttt{cap\_storage}  & output storage capacitor bank     & $213.5\times119.4$ & 19.998\,pF & \PASS \\
\bottomrule
\end{longtable}

Transistor modules were checked with foundry DRC + \texttt{netgen} LVS. Capacitor modules were
checked with foundry DRC + a plate-separation test, because no extractor available here models
a MIM device.

\section{The MIM-B capacitor module}
\label{sec:mim}

The previous design used GLayout's \texttt{mimcap()}, which draws Metal2 $+$ \texttt{CAP\_MK}
$+$ a Via2 array $+$ Metal3. On GF180 that is not a capacitor: it is an ordinary via array, a
short circuit. A real MIM needs \textbf{FuseTop (75/0)} as the top electrode ---
\texttt{mim\_a.drc} states \texttt{mim7\_l1 = fusetop.not(cap\_mk)}, i.e.\ \texttt{CAP\_MK}
only \emph{marks} FuseTop.

Which metal pair carries the MIM depends on the metal stack. The deck's variant table is
authoritative:

\begin{center}
\begin{tabular}{@{}ccccl@{}}
\toprule
\textbf{Variant} & \textbf{metal\_top} & \textbf{mim\_option} & \textbf{metal\_level} & \\
\midrule
A & 30K & A & 3LM & \\
B & 11K & B & 4LM & \\
C & 9K  & \textbf{B} & \textbf{5LM} & $\leftarrow$ candidate \\
D & 11K & \textbf{B} & \textbf{5LM} & $\leftarrow$ used for sign-off \\
E & 9K  & B & 6LM & \\
F & 9K  & A & 6LM & \\
\bottomrule
\end{tabular}
\end{center}

This design is \textbf{5LM} (met1--met5; the padring cell \texttt{gf180mcu\_fd\_io\_\_asig\_5p0}
is a 5LM cell), so only \texttt{mim\_option=B} exists. The module is therefore built as:
Metal4 bottom plate, FuseTop top plate, Via4 contact, and \emph{both} terminals brought out on
Metal5 --- rule MIMTM.10 forbids contacting the Metal4 bottom plate from below.

\subsection*{MIM-B rule compliance}

\begin{longtable}{@{}llll@{}}
\toprule
\textbf{Rule} & \textbf{Requirement} & \textbf{Implemented} & \textbf{Result} \\
\midrule
\endhead
MIMTM.1  & met4 to neighbouring MIM bottom plate $\geq 1.2$\um & 1.4\um  & \PASS \\
MIMTM.2  & met4 encloses via4 in the MIM region $\geq 0.4$\um  & 0.5\um  & \PASS \\
MIMTM.3  & met4 encloses FuseTop $\geq 0.6$\um                 & 0.8\um  & \PASS \\
MIMTM.4  & FuseTop encloses via4 $\geq 0.4$\um                 & 0.4\um  & \PASS \\
MIMTM.5  & FuseTop to bottom-contact via4 $\geq 0.4$\um        & 0.5\um  & \PASS \\
MIMTM.6  & FuseTop to unrelated FuseTop $\geq 0.6$\um          & 1.4\um  & \PASS \\
MIMTM.7  & FuseTop covered by \texttt{CAP\_MK} $\geq 0$\um     & 0.2\um  & \PASS \\
MIMTM.8a & FuseTop area $\geq 25$\umsq                         & enforced in code & \PASS \\
MIMTM.8b & FuseTop area $\leq 10000$\umsq\ per capacitor       & enforced in code & \PASS \\
MIMTM.9  & via4 spacing on the top plate $\geq 0.5$\um         & 0.54\um & \PASS \\
MIMTM.10 & no Via3 may touch the Metal4 bottom plate           & both terminals on met5 & \PASS \\
MIMTM.11 & FuseTop per shared bottom plate $\leq 10000$\umsq   & one plate per unit & \PASS \\
\bottomrule
\end{longtable}

\textbf{Critical caveat.} A clean DRC result is only meaningful if the MIM rules actually ran.
Every MIM rule keys off \texttt{fusetop}; with no FuseTop drawn they match nothing and pass
\emph{vacuously} --- which is exactly why the previous, shorted design was reported ``DRC
clean''. The sign-off log must therefore be read, not just its verdict. For the submitted file
it contains:

\begin{center}
\texttt{MIM Option selected: B} \qquad \texttt{fusetop has 34 polygons}
\end{center}

\section{Metal, via and density rules}

\begin{longtable}{@{}lllll@{}}
\toprule
\textbf{Rule} & \textbf{Requirement} & \textbf{Design uses} & \textbf{Measured} & \textbf{Result} \\
\midrule
\endhead
M1.1 & met1 width $\geq 0.23$\um  & 0.5\um straps         & min 0.500\um & \PASS \\
M2.1 & met2 width $\geq 0.28$\um  & 8.0\um pin landings   & min 0.500\um & \PASS \\
M3.1 & met3 width $\geq 0.28$\um  & 0.8/2.0\um lanes      & min 0.500\um & \PASS \\
M4.1 & met4 width $\geq 0.28$\um  & 2.0\um rails          & min 0.500\um & \PASS \\
M5.1 & met5 width $\geq 0.28$\um  & 2.0\um buses          & min 2.000\um & \PASS \\
V4.1 & via4 exactly $0.26$\um     & 0.26\um               & 0.26\um      & \PASS \\
V4.2b& via4 spacing in arrays $\geq0.36$\um & 0.54\um pitch gap & ---   & \PASS \\
M1.4 & met1 coverage $> 30\%$     & dummy fill            & \textbf{43.2\%} & \PASS \\
M2.4 & met2 coverage $> 30\%$     & dummy fill            & \textbf{43.7\%} & \PASS \\
M3.4 & met3 coverage $> 30\%$     & dummy fill            & \textbf{42.8\%} & \PASS \\
M4.4 & met4 coverage $> 30\%$     & dummy fill            & \textbf{45.6\%} & \PASS \\
M5.4 & met5 coverage $> 30\%$     & dummy fill            & \textbf{45.9\%} & \PASS \\
PL.8 & poly2 coverage $= 14\%$    & dummy fill            & \textbf{43.2\%} & \PASS \\
\bottomrule
\end{longtable}

Dummy fill uses a $1.4$\um\ clearance on met4 rather than the $0.3$\um\ metal spacing, because
MIMTM.1 keeps any Metal4 at least $1.2$\um\ away from a MIM bottom plate. Fill also keeps clear
of the slot edge, the Metal2 corner keep-outs and the five pin landings.

\section{Project-slot rules}

\begin{longtable}{@{}llll@{}}
\toprule
\textbf{Requirement} & \textbf{Source} & \textbf{Measured} & \textbf{Result} \\
\midrule
\endhead
Area exactly $550\times550$\um & \texttt{D05\_EH.def} \texttt{DIEAREA} & bbox $(0,0)$--$(550,550)$ & \PASS \\
Layer 0/0 boundary present & organiser flow & present & \PASS \\
Five pins, Metal2, west edge & \texttt{D05\_EH\_interface.yaml} & all five landed & \PASS \\
\quad VSS (W12, \texttt{dvss})   & 6 rects, 60.000\umsq & 60.000\umsq covered & \PASS \\
\quad RFP (W13, \texttt{asig\_5p0}) & 8 rects, 20.320\umsq & 20.320\umsq covered & \PASS \\
\quad RFN (W14, \texttt{asig\_5p0}) & 8 rects, 20.320\umsq & 20.320\umsq covered & \PASS \\
\quad VOUT (W15, \texttt{asig\_5p0}) & 8 rects, 20.320\umsq & 20.320\umsq covered & \PASS \\
\quad VCC (W16, \texttt{dvdd})   & 6 rects, 60.000\umsq & 60.000\umsq covered & \PASS \\
Metal2 corner keep-outs & \texttt{D05\_EH.def} \texttt{BLOCKAGES} & none issued for EH & \PASS \\
\quad (regression against variant D) & $(529,548)$--$(550,550)$, $(0,0)$--$(2,21)$ & both clear on a D build & \PASS \\
Hierarchical GDS, unique cell names & requested by the integrator & 163 cells, all \texttt{D05\_}-prefixed, no \texttt{\$} & \PASS \\
\bottomrule
\end{longtable}

\section{Verification methodology}
\label{sec:method}

Three points matter for interpreting the results.

\textbf{DRC is the foundry KLayout deck, not magic.} magic's \texttt{gf180mcuD} technology has
no MIM device and a much thinner rule set; it is what allowed the previous shorted capacitors
to pass. All DRC numbers here come from
\texttt{run\_drc.py -{-}variant=D -{-}density -{-}antenna}.

\textbf{LVS runs on the capacitor-free FET view.} No extractor available here models a MIM, so
a with-capacitor extraction would read the via4 sea on FuseTop as a plate short. LVS therefore
proves the transistor topology; the capacitors are verified geometrically instead.

\textbf{Capacitors are verified by layer-aware net extraction.} A via4 landing \emph{on}
FuseTop is modelled as met5$\leftrightarrow$FuseTop; a via4 landing \emph{off} FuseTop as
met5$\leftrightarrow$met4; FuseTop is never joined to met4 --- that gap \emph{is} the
capacitor. A naive ``via4 joins met4 to met5'' model reports every MIM as a short.

\section{Verification results}

\begin{longtable}{@{}lll@{}}
\toprule
\textbf{Check} & \textbf{Result} & \textbf{Status} \\
\midrule
\endhead
Foundry DRC, geometry $+$ density $+$ antenna & 0 flagged items & \PASS \\
\quad MIM rules actually executed & \texttt{MIM Option: B}, \texttt{fusetop has 34 polygons} & \PASS \\
\texttt{netgen} LVS, exact file and top cell & circuits match uniquely, 9 ports & \PASS \\
Per-module DRC + LVS (9 modules) & all clean & \PASS \\
Net separation & 9 named nets, each on exactly one net & \PASS \\
Dangling metal & 0.00\umsq\ unreachable & \PASS \\
Capacitor plate separation & 34 plates, none shorted or floating & \PASS \\
Capacitor net assignment & all 5 on their intended net pairs & \PASS \\
Route widths & nothing below its DRC minimum & \PASS \\
PEX device parameters & schematic histogram $=$ extracted histogram & \PASS \\
\bottomrule
\end{longtable}

\subsection*{Capacitor net assignment}

\begin{center}
\begin{tabular}{@{}llll@{}}
\toprule
\textbf{Bank} & \textbf{Top plate} & \textbf{Bottom plate} & \textbf{Intended} \\
\midrule
storage (28 units, 19.998\,pF) & VOUT  & VSS & VSS / VOUT \quad\PASS \\
stage 1 (2 units, 1.200\,pF)   & N1    & RFN & RFN / N1 \quad\PASS \\
stage 2 (2 units, 1.200\,pF)   & N2    & RFP & RFP / N2 \quad\PASS \\
inject 1 (1 unit, 0.300\,pF)   & VMID1 & RFP & RFP / VMID1 \quad\PASS \\
inject 2 (1 unit, 0.300\,pF)   & VMID2 & RFN & RFN / VMID2 \quad\PASS \\
\bottomrule
\end{tabular}
\end{center}

\section{Measured outputs}

Post-layout, on the extracted netlist (2293 R, 439 C) with the capacitors at the values the
MIM-B geometry actually realises, simulated in \texttt{ngspice-46}.

\begin{center}
\begin{tabular}{@{}lll@{}}
\toprule
\textbf{Quantity} & \textbf{Condition} & \textbf{Result} \\
\midrule
VRECT (rectified rail) & $\pm0.9$\,V, 30\,MHz differential & 0.817\,V \\
VOUT (harvested output) & same & 1.578\,V \\
Pump gain VOUT/VRECT & same & $2.46\times$ \\
Input impedance $Z_\mathrm{in}$ & 30\,MHz, large-signal & $2324 - 170j\ \Omega$ differential \\
Matching network (off-chip) & to 100\,$\Omega$ & shunt 10.5\,pF $+$ series 1265\,nH/side \\
VOUT at 0\,dBm & 30\,MHz, unmatched & 0.31\,V \\
VOUT at $+8$\,dBm & 30\,MHz, unmatched & 1.94\,V \\
Sensitivity (VOUT $\geq 1$\,V) & unmatched & $+3.7$\,dBm \\
Sensitivity (VOUT $\geq 1$\,V) & L-matched & $-7.2$\,dBm \\
Matching gain & --- & $\approx 10.8$\,dB \\
Total on-chip MIM capacitance & --- & 22.998\,pF \\
\bottomrule
\end{tabular}
\end{center}

\section{Defects found and corrected}

\begin{longtable}{@{}p{0.30\textwidth}p{0.42\textwidth}l@{}}
\toprule
\textbf{Defect} & \textbf{Cause and correction} & \textbf{Status} \\
\midrule
\endhead
Capacitors were metal shorts & \texttt{mimcap()} drew met2/via2/met3 with no FuseTop; the 22\,pF storage bank tied VOUT to VSS. Rebuilt on the MIM-B stack. & fixed \\
Wrong metal pair & \texttt{mim\_option=A} (met2/met3) exists only for 3LM/6LM; this is 5LM. Capacitors moved to met4/met5. & fixed \\
Flat GDS, cell named \texttt{rf\_energy\_harvester\$1} & \texttt{component\_snap\_to\_grid()} flattens by design. Skipped; cells renamed and prefixed. & fixed \\
No connection to the padring & The core's rail pads sat 137\um\ inside the slot on met4. Added a met4/met3 boundary router to the Metal2 landings. & fixed \\
Density 0.1--11\,\% & Added dummy fill on all six layers. & fixed \\
Five met4 rails welded together & \texttt{mlink} draws each stub on the \emph{port's own} layer; met4 ports shorted the rails they crossed. Ports now drop on met3. & fixed \\
N1/N2 shorted to VSS & Those ports sit \emph{inside} the pump, above its own met3 VSS bar. They now route upward to a second rail band. & fixed \\
N1 merged with VMID1 & Two capacitors shared an $x$ column, so their met5 buses overlapped. Each capacitor now has its own column. & fixed \\
VOUT and VRECT each split in two & The bottom rail did not extend to the met5 tie column, so the tie landed on an isolated pad. Rails now include the tie anchor. & fixed \\
\bottomrule
\end{longtable}

\section{Open items}

\begin{enumerate}
\item \OPEN\ \textbf{Process variant.} The capacitors are MIM-B because that is what the deck
offers for a 5LM stack. The organisers should confirm the shuttle runs a 5LM variant with the
MIM module enabled. A different metal count moves the capacitors to a different metal pair, and
the floorplan with them.
\item \OPEN\ \textbf{Post-layout re-closure.} The results in Sec.~6 use the MIM-B capacitance
values but an extraction taken from the capacitor-free view. They have not yet been re-closed
against the real MIM-B parasitics. This is the remaining engineering step before tapeout.
\end{enumerate}

\end{document}
